\chapter{Introducción}
\label{chap:introduction}

Este capítulo establece el contexto, la motivación y el problema de investigación que aborda la presente tesis, cuyo enfoque central radica en el diseño de una arquitectura de razonamiento estructural, basada en \ac{RI}, para mitigar las alucinaciones en la generación \emph{text-to-query} de consultas empresariales complejas, específicamente orientada al estándar híbrido \ac{SQL}/\ac{PGQ}.

Esta investigación se sitúa en la intersección entre los \acp{LLM}, los sistemas de gestión de bases de datos relacionales (\acp{RDBMS}) y los lenguajes de consulta híbridos, pues estos lenguajes combinan el paradigma relacional tradicional con las consultas sobre grafos de propiedades (\emph{property graphs}), un modelo recientemente formalizado en estándares internacionales como ISO/IEC \cite{ISO9075_16_2023}.

Como se evidenciará en el posterior desarrollo, el diseño propuesto representa una respuesta estructurada al creciente interés de la industria por desplegar interfaces en lenguaje natural altamente fiables en entornos empresariales a gran escala. Reducir las alucinaciones y mantener la consistencia lógica constituyen retos críticos actuales \cite{WrenAI2025, TA-SQL2024}, los cuales demandan métodos avanzados que integren el razonamiento tabular y de grafos \cite{STRuCT-LLM2026}.


%A pesar de que la generación text-to-queries (Text→Query) con modelos de lenguaje grandes \ac{LLM} se ha desarrollado rápidamente, en situaciones empresariales todavía presenta fallas sistemáticas en términos de seguridad operacional, fidelidad al esquema y corrección semántica. Estas fallas se presentan en forma de alucinaciones específicas de consulta: referencias a columnas/tablas/atributos que no existen, rutas de uniones inventadas, errores en las agregaciones y malentendidos en las condiciones. En el caso de los grafos, son patrones o etiquetas de recorrido que no concuerdan con el modelo gráfico. Investigaciones actuales sobre evaluación "realista" indican que el desempeño en benchmarks tradicionales no se refleja correctamente en los flujos empresariales: 

\section{Motivación y Contexto}
\label{sec:motivacion_contexto}

En el contexto actual de la adopción de \acp{LLM} en procesos empresariales y el creciente interés por integrarlos en la optimización de flujos de trabajo, se evidencia una clara tendencia hacia interfaces basadas en lenguaje natural \cite{Data2Decisions}, en detrimento de las interfaces gráficas tradicionales. Este cambio es particularmente notable en organizaciones \emph{data-driven}, cuya toma de decisiones depende de la capacidad de responder preguntas de negocio en ciclos cortos y con altas garantías de calidad \cite{llm-insight}. En este marco, las interfaces en lenguaje natural para la consulta de datos han resurgido con fuerza, habilitando experiencias de \emph{self-service analytics}. Estas permiten que perfiles no técnicos formulen preguntas en texto libre y lenguaje cotidiano, obteniendo resultados estructurados sin la necesidad de escribir manualmente consultas complejas \cite{NLinterfaces}. 

Sin embargo, el despliegue de estos modelos en entornos de producción demanda el cumplimiento de tres pilares críticos: (i) \emph{integridad técnica}, que garantiza tanto la ejecutabilidad en el \ac{RDBMS} como la fidelidad semántica de la consulta \cite{spider2_2025}; (ii) \emph{seguridad y gobernanza}, mediante mecanismos de control de acceso basado en roles (RBAC) y auditoría para prevenir la exfiltración de datos sensibles \cite{SAFENLIDB}; y (iii) \emph{robustez y escalabilidad del esquema}, asegurando que el modelo mantenga su precisión frente a catálogos de datos masivos, dinámicos y con nomenclaturas complejas \cite{llmfk}.


Asimismo, a pesar de los avances acelerados, los \acp{LLM} aún presentan fallas sistemáticas de \emph{fidelidad y grounding} al generar artefactos formales. Específicamente en la tarea de \emph{text-to-query}, estas fallas se manifiestan como referencias a tablas o columnas inexistentes, rutas de \emph{join} inventadas, condiciones lógicas mal interpretadas o agregaciones incongruentes con el esquema. Si bien esta es una limitación general de los modelos de lenguaje, el riesgo se agrava severamente cuando el código se genera para un \ac{RDBMS}: un error no solo afecta la exactitud analítica, sino que puede inducir sobrecostos computacionales, vulnerabilidades de seguridad e incluso modificar datos existentes. Investigaciones recientes demuestran que el alto desempeño en \emph{benchmarks} clásicos decae drásticamente al enfrentarse a bases de datos corporativas reales, ruidosas y con flujos de trabajo de múltiples pasos \cite{li2023bird, spider2_2025}, lo cual refuerza la necesidad de transitar desde enfoques de \emph{prompting} monolíticos y arquitecturas puramente agénticas hacia sistemas más robustos y orquestados \cite{EvaluatingRobustness}.


A su vez en paralelo, el ecosistema de bases de datos está evolucionando hacia modelos híbridos. En dominios como prevención de fraude, ciberseguridad, logística y \emph{compliance}, el modelado de relaciones mediante grafos resulta más natural para expresar caminos y vecindarios complejos. No obstante, dado que las bases de datos relacionales siguen dominando el entorno corporativo por su estabilidad, la integración de \emph{property graphs} se presenta como una evolución natural de la industria. Esto se valida con la reciente estandarización ISO/IEC 9075-16:2023 (\ac{SQL}/\ac{PGQ}), la cual formaliza la consulta de grafos de propiedades directamente desde \ac{SQL} \cite{ISO9075_16_2023}. Si bien esto unifica ambos paradigmas, también incrementa significativamente la dificultad de la tarea \emph{text-to-query}: el espacio de búsqueda se expande y la validación lógica se vuelve más rigurosa, multiplicando el riesgo de alucinaciones.

En la industria actual, las soluciones comerciales suelen apoyarse en estrategias heurísticas como inyectar el esquema en el \emph{prompt} o aplicar ciclos de autocorrección mediante retroalimentación en texto libre. Estos métodos han demostrado ser frágiles ante la ambigüedad semántica del usuario y el \emph{schema drift} \cite{AmbiSQL}. Por su parte, la literatura reciente aborda este reto desde diversos frentes: decodificación restringida por gramática para garantizar validez sintáctica \cite{PICARD2022}, arquitecturas modulares que dividen la carga de razonamiento \cite{DIN-SQL2023}, técnicas de alineación pre-generación para mitigar alucinaciones \cite{TA-SQL2024}, métodos de evaluación sin \emph{ground truth} como el \emph{metamorphic testing} \cite{WrenAI2025}, y la unificación del razonamiento tabular y de grafos \cite{STRuCT-LLM2026}, entre otros. 

Es así que estos antecedentes convergen en la premisa central de este trabajo: en entornos empresariales, la mitigación efectiva de alucinaciones sobre dialectos híbridos requiere un diseño arquitectónico \emph{sistémico}. Por ello, la presente investigación se motiva por la necesidad de diseñar una arquitectura que:
\begin{itemize}
    \item Imponga un \emph{anclaje estricto} (\emph{schema grounding}) a nivel relacional y de grafo para erradicar las referencias inventadas \cite{SchemaGraphSQL}.
    \item Utilice una \ac{RI} tipada que actúe como contrato determinista entre la comprensión del lenguaje natural y la síntesis del código final \cite{atscale2024sml}.
    \item Incorpore mecanismos de \emph{verificación activa} (estática y por ejecución controlada) antes de la delegación en producción \cite{CodeScore}.
    \item Soporte nativamente el estándar \ac{SQL}/\ac{PGQ}, garantizando las exigencias empresariales de rendimiento, seguridad y reproducibilidad.
    \item Asegure la transferibilidad tecnológica mediante el uso de herramientas de orquestación y motores de base de datos compatibles con el stack empresarial contemporáneo.
\end{itemize}


\section{Planteamiento del Problema}

La adopción de \acp{LLM} en entornos empresariales ha dado lugar a propuestas de interfaces en lenguaje natural que frecuentemente carecen de fundamentos técnicos sólidos, teniendo como resultado que en lugar de traducir consultas con fidelidad semántica, se generan soluciones mal fundamentadas que incurren en sobrecostos operacionales y exhiben escasa adaptabilidad a cambios y escala de proyectos.

Esta brecha se origina en que los \acp{LLM} aplicados a \emph{text-to-query} tienden a generar consultas \ac{SQL}/\ac{PGQ} plausibles pero incorrectas, debido a fallas de anclaje al esquema, ambigüedad semántica o ausencia de restricciones durante la generación. En entornos empresariales, estas fallas se manifiestan como \emph{alucinaciones} que comprometen la corrección, la seguridad y el rendimiento. El problema se agudiza en el estándar híbrido \ac{SQL}/\ac{PGQ}, el cual, como se explicó en la sección anterior, amplía considerablemente el espacio de expresividad de las consultas, pero también incrementa la dificultad de garantizar su corrección.

Por tanto, el problema que aborda esta tesis es: \textit{¿cómo diseñar e implementar una arquitectura de razonamiento estructural basada en \ac{RI} que mitigue sistemáticamente las alucinaciones en la generación \emph{text-to-query} y produzca consultas \ac{SQL}/\ac{PGQ} correctas, verificables y operables bajo restricciones empresariales, con garantías de escalabilidad y transferibilidad entre proyectos?}


\section{Objetivos}

\textbf{Objetivo general:} Diseñar, implementar y evaluar experimentalmente una arquitectura para la generación \emph{text-to-query} orientada al estándar \ac{SQL}/\ac{PGQ}, basada en razonamiento estructural y una \ac{RI} tipada, capaz de reducir alucinaciones de esquema y errores semánticos en escenarios representativos de consulta empresarial, y cuya eficacia se contraste con baselines del estado del arte mediante métricas de corrección, alucinación, seguridad y eficiencia sobre benchmarks establecidos como Spider \cite{spider2_2025}, BIRD \cite{li2023bird} y Text2GQL-Bench \cite{text2gqlbench}.

\subsection{Objetivos Específicos}

\begin{itemize}
    \item Caracterizar el fenómeno de alucinación en \emph{text-to-query} (relacional e híbrido) mediante una taxonomía operacional de errores (p.\,ej., referencias inexistentes, rutas inventadas, condiciones/agrupaciones incorrectas), y mapear sus causas principales en \acp{LLM} bajo \emph{schema drift} y ambigüedad.
    \item Diseñar una \ac{RI} unificada para consultas híbridas que represente componentes relacionales (proyección, selección, \emph{joins}, agregación) y componentes de patrón de grafo (variables, aristas, restricciones), junto con reglas de tipado, canoni\-calización y compilación determinista conforme al estándar \ac{SQL}/\ac{PGQ} \cite{ISO9075_16_2023}.
    \item Construir un \emph{pipeline} de razonamiento estructural que integre: anclaje al esquema, generación restringida (p.\,ej., decodificación guiada por gramática o validación incremental en el espíritu de \cite{PICARD2022}), y un bucle de verificación que combine chequeos estáticos y ejecución controlada (sandbox) con criterios de costo/seguridad.
    \item Definir un protocolo experimental y un conjunto de evaluación que cubra \emph{text-to-SQL} y \emph{text-to-graph queries} (incluyendo benchmarks enterprise y de grafos), y evaluar la propuesta frente a métodos comparables (p.\,ej., \cite{DIN-SQL2023,TA-SQL2024,CRED-SQL2025}) usando métricas de exactitud por ejecución, tasa de alucinación de esquema, robustez a drift y sobrecosto de latencia.
\end{itemize}

\section{Organización de la tesis}
El resto de este documento está organizado de la siguiente manera:

\begin{itemize}
	\item El capítulo~\ref{chap:background} describe el marco teórico, para un mejor entendimiento del alcance de \acp{LLM} y alucinaciones, fundamentos de \emph{text-to-SQL}, bases de \ac{SQL}/\ac{PGQ}, y técnicas de \ac{RI}, decodificación restringida y verificación.
	\item El capítulo~\ref{chap:stateoftheart} revisa el estado del arte en arquitecturas \emph{text-to-query} (pipelines por etapas, agentes, \emph{constrained decoding}, detección/mitigación de alucinaciones) y posiciona el problema en el contexto empresarial y del estándar híbrido.
	\item El capítulo~\ref{chap:proposal} describe la propuesta: la \ac{RI} tipada para \ac{SQL}/\ac{PGQ}, el compilador \ac{RI} $\rightarrow$ \ac{SQL}/\ac{PGQ}, el pipeline de razonamiento estructural y el bucle de verificación (estático y por ejecución).
	\item El capítulo~\ref{chap:experiments} detalla el diseño experimental: datasets/benchmarks, construcción de corpus híbrido para \ac{SQL}/\ac{PGQ}, métricas de alucinación y corrección, baselines y \emph{ablations}.
	\item El capítulo~\ref{chap:results} presenta y analiza los resultados obtenidos, incluyendo robustez a \emph{schema drift}, seguridad/eficiencia, y evaluación de componentes, validando los resultados con datasets relevantes.
	\item El capítulo~\ref{chap:conclusions} concluye la tesis, resume aportes, limitaciones y líneas de trabajo futuro.
\end{itemize}



\section{Cronograma}
\definecolor{headerblue}{RGB}{51,102,153}

{\footnotesize
\setlength{\tabcolsep}{3pt}
\setlength{\LTleft}{0pt}
\setlength{\LTright}{0pt}

\begin{longtable}{|>{\centering\arraybackslash}p{0.45cm}|>{\centering\arraybackslash}p{0.45cm}|>{\centering\arraybackslash}p{0.4cm}|p{3.2cm}|p{5.6cm}|>{\centering\arraybackslash}p{1.2cm}|>{\centering\arraybackslash}p{1.7cm}|>{\centering\arraybackslash}p{0.45cm}|}

\caption{Cronograma de actividades de la tesis.} \label{tab:cronograma} \\

\hline
\rowcolor{headerblue}
 &
 &
\textcolor{white}{\textbf{N.°}} &
\textcolor{white}{\textbf{Actividad}} &
\textcolor{white}{\textbf{Detalles}} &
\textcolor{white}{\textbf{Semana}} &
\textcolor{white}{\textbf{Fecha}} &
\textcolor{white}{\textbf{\%}} \\
\hline
\endfirsthead

\hline
\rowcolor{headerblue}
 &
 &
\textcolor{white}{\textbf{N.°}} &
\textcolor{white}{\textbf{Actividad}} &
\textcolor{white}{\textbf{Detalles}} &
\textcolor{white}{\textbf{Semana}} &
\textcolor{white}{\textbf{Fecha}} &
\textcolor{white}{\textbf{\%}} \\
\hline
\endhead

\hline
\multicolumn{8}{r}{\textit{Continúa en la siguiente página...}} \\
\endfoot

\hline
\endlastfoot

% ==============================
% SEMESTRE 2026-1
% ==============================
\multirow{19}{*}{\rotatebox[origin=c]{90}{\textbf{Semestre 2026-1}}} &
\multirow{12}{*}{\rotatebox[origin=c]{90}{Parcial}} &
1 & Alineamiento inicial del trabajo &
Revisión del curso, entregables, estructura de la tesis y ajuste del tema con enfoque en SQL/PGQ + representación intermedia. &
Sem.\ 1 & 11/03 -- 14/03 & 2 \\
\cline{3-8}

& & 2 & Definición final de tema y asesor &
Alcance, contribución central y validación con asesora. Primera revisión de tabla survey. &
Sem.\ 2 & 14/03 -- 21/03 & 5 \\
\cline{3-8}

& & 3 & Formalización y delimitación &
Compromiso de asesoría, planteamiento del problema, objetivo general y específicos. &
Sem.\ 3 & 21/03 -- 28/03 & 10 \\
\cline{3-8}

& & 4 & Consolidación del survey &
Depuración del estado del arte, clasificación de trabajos, taxonomía y bibliografía base. &
Sem.\ 4 & 28/03 -- 01/04 & 15 \\
\cline{3-8}

& & 5 & Redacción de Introducción v1 &
Motivación, contexto, problema, objetivos, organización y cronograma. &
Sem.\ 5 & 01/04 -- 08/04 & 18 \\
\cline{3-8}

& & 6 & Cierre de Introducción &
Ajustes con feedback y versión entregable del capítulo Introducción. &
Sem.\ 5 & 08/04 -- 11/04 & 20 \\
\cline{3-8}

& & 7 & Redacción de Marco Teórico &
Fundamentos teóricos, conceptos clave y modelos explicativos de la investigación. &
Sem.\ 6 & 11/04 -- 14/04 & 28 \\
\cline{3-8}

& & 8 & Redacción de Estado del Arte &
Literatura reciente, trabajos previos, vacíos de investigación y posicionamiento. &
Sem.\ 6 & 14/04 -- 18/04 & 35 \\
\cline{3-8}

& & 9 & Capítulo Propuesta v1 &
Arquitectura, pipeline, representación intermedia, componentes y flujo de verificación. &
Sem.\ 7 & 18/04 -- 25/04 & 45 \\
\cline{3-8}

& & 10 & Capítulo Propuesta v2 &
Protocolo experimental, datasets, métricas, baselines y criterios de evaluación. &
Sem.\ 8 & 26/04 -- 02/05 & 50 \\
\cline{3-8}

& & 11 & Integración del Plan de Tesis &
Consolidación de capítulos, coherencia interna, flujo narrativo y bibliografía final. &
Sem.\ 9 & 03/05 -- 09/05 & 55 \\
\cline{3-8}

& & 12 & Plan de Tesis listo para envío &
Cierre de documento, diapositivas y ensayo previo a la exposición. &
Sem.\ 10 & 11/05 -- 16/05 & 55 \\
\cline{2-8}

& \multirow{7}{*}{\rotatebox[origin=c]{90}{Final}} &
13 & Implementación base experimental &
Repositorio, estructura del prototipo, casos de prueba y baseline(s). &
Sem.\ 11 & 17/05 -- 23/05 & 62 \\
\cline{3-8}

& & 14 & Implementación arquitectura v1 &
Flujo NL $\rightarrow$ RI $\rightarrow$ consulta estructurada/verificable. &
Sem.\ 12 & 24/05 -- 30/05 & 68 \\
\cline{3-8}

& & 15 & Integración end-to-end y validación &
Integración de componentes, validaciones estructurales y depuración. &
Sem.\ 13 & 31/05 -- 06/06 & 73 \\
\cline{3-8}

& & 16 & Experimentación preliminar &
Pruebas comparativas, recopilación de resultados y análisis preliminar. &
Sem.\ 14 & 07/06 -- 13/06 & 77 \\
\cline{3-8}

& & 17 & Resultados al 80 y Avances &
Prototipo funcional, resultados parciales sólidos y figuras iniciales. &
Sem.\ 15 & 14/06 -- 20/06 & 80 \\
\cline{3-8}

& & 18 & Cierre del documento de Avances &
Documento completo, conclusiones preliminares y preparación de diapositivas. &
Sem.\ 16 & 21/06 -- 27/06 & 80 \\
\cline{3-8}

& & 19 & Envío y exposición de Avances &
Presentación final del semestre con resultados parciales consolidados. &
Sem.\ 17 & 28/06 -- 04/07 & 80 \\
\hline

% ==============================
% PERÍODO INTERMEDIO
% ==============================
\multicolumn{2}{|c|}{\textit{Int.}} &
20 & Revisión y fortalecimiento &
Análisis de trabajos complementarios y refinamiento de la arquitectura propuesta. &
 & 10/07 -- 10/08 & 82 \\
\hline

% ==============================
% SEMESTRE 2026-2
% ==============================
\multirow{9}{*}{\rotatebox[origin=c]{90}{\textbf{Semestre 2026-2}}} &
\multirow{4}{*}{\rotatebox[origin=c]{90}{Parcial}} &
21 & Proyecto de Investigación &
Propuesta formal: impacto, viabilidad y posible aplicación institucional. &
Sem.\ 1--2 & 12/08 -- 22/08 & 84 \\
\cline{3-8}

& & 22 & Ampliación experimental &
Nuevos escenarios, ablations, análisis de errores y validación empresarial. &
Sem.\ 3--5 & 23/08 -- 15/09 & 87 \\
\cline{3-8}

& & 23 & Resultados finales y discusión &
Comparación con baselines, análisis crítico y discusión de aportes. &
Sem.\ 6--8 & 16/09 -- 06/10 & 90 \\
\cline{3-8}

& & 24 & Preparación y Exposición Parcial &
Presentación con metodología, resultados preliminares y validación ante jurado. &
Sem.\ 9 & 07/10 -- 13/10 & 91 \\
\cline{2-8}

& \multirow{5}{*}{\rotatebox[origin=c]{90}{Final}} &
25 & Revisión y ajustes post-parcial &
Correcciones metodológicas y refinamiento con feedback del jurado. &
Sem.\ 10 & 14/10 -- 20/10 & 92 \\
\cline{3-8}

& & 26 & Síntesis de resultados y divulgación &
Hallazgos en formato de reporte técnico o borrador de artículo. &
Sem.\ 11 & 21/10 -- 27/10 & 93 \\
\cline{3-8}

& & 27 & Redacción final de tesis &
Cierre de capítulos, conclusiones, trabajo futuro y referencias. &
Sem.\ 12--14 & 28/10 -- 15/11 & 96 \\
\cline{3-8}

& & 28 & Correcciones con asesora &
Ajustes de fondo y forma, normalización y aprobación final. &
Sem.\ 15--16 & 16/11 -- 30/11 & 98 \\
\cline{3-8}

& & 29 & Preparación y Exposición Final &
Diapositivas, ensayo de tiempos y sustentación final ante jurado. &
Sem.\ 17--18 & 01/12 -- 12/12 & 100 \\

\end{longtable}
}

